\section*{Capítulo \ref{ch:b3:mantle-convection} --- Convección del manto y fuerzas motoras de las placas}

\begin{solution}{exo:b3:mantle-convection:1}
$\mathrm{Ra} = 3500\times9.9\times\num{3e-5}\times1000\times(\num{6.7e5})^3/
(10^{-6}\times\num{5e20}) = \num{6.25e5}$, unas 950 veces el \omterm{def:b3:mantle-convection:critical}{número de Rayleigh crítico}
657.5: la capa convecta con vigor.
\end{solution}

\begin{solution}{exo:b3:mantle-convection:2}
$\tau_M = 10^{21}/\num{7.7e10} = \qty{1.30e10}{s}$, unos \qty{412}{yr}. La
marea, que dura horas, encuentra un manto elástico; el deshielo de un casquete de hielo,
durante miles de años, y el hundimiento de una losa, durante millones, encuentran un
manto viscoso.
\end{solution}

\begin{solution}{exo:b3:mantle-convection:3}
$\tau = 4\pi\times10^{21}/(3300\times9.8\lambda)$: \qty{3.89e11}{s} =
\qty{12.3}{kyr} para \qty{1000}{km}, y \qty{4.1}{kyr} para \qty{3000}{km}. La
carga más extensa se recupera antes.
\end{solution}

\begin{solution}{exo:b3:mantle-convection:4}
$g\rho_m\alpha(T_m - T_0)\kappa t = 9.8\times3300\times\num{3.1e-5}\times1450
\times\num{0.8e-6}\times50\times\num{3.156e13} = \qty{1.835e12}{N.m^{-1}}$; el
corchete es $1 + 2\times3300\times\num{3.1e-5}\times1450/(\pi\times2270) = 1.0416$,
así que $F_{rp} = \qty{1.91e12}{N.m^{-1}}$, 24 veces menor que el tirón de la losa.
\end{solution}

\begin{solution}{exo:b3:mantle-convection:5}
(a) $v = 3300\times9.8\times\num{3e-5}\times300\times(\num{e4})^2/10^{21} =
\qty{2.91e-11}{m.s^{-1}}$; el tiempo de ascenso a lo largo de \qty{10}{km} es $d/v =
\qty{3.44e14}{s}$, \qty{10.9}{Myr}, el tiempo de difusión $d^2/\kappa =
\qty{e14}{s}$, \qty{3.2}{Myr}. Su razón es $(d^2/\kappa)/(d/v) =
vd/\kappa = \rho g\alpha\Delta Td^3/(\kappa\eta) = \mathrm{Ra} = 0.29$. (b) Se
enfría primero: las masas pequeñas no convectan. (c) $\mathrm{Ra} \propto d^3$, así que $d =
10\times(657.5/0.291)^{1/3} = \qty{131}{km}$.
\end{solution}

\begin{solution}{exo:b3:mantle-convection:6}
$2\mathrm{Ra}_c = 1315$ y $\mathrm{Ra}/2\mathrm{Ra}_c = 7604$, cuya raíz cúbica
es 19.66. $\delta = 2891/19.66 = \qty{147}{km}$, $\mathrm{Nu} = 19.66/2 =
9.8$, $u = (10^{-6}/\num{2.891e6})\times19.66^2 = \qty{1.34e-10}{m.s^{-1}} =
\qty{0.42}{cm/yr}$, y $q = 9.8\times3.1\times1450/\num{2.891e6} =
\qty{15}{mW.m^{-2}}$.
\end{solution}

\begin{solution}{exo:b3:mantle-convection:7}
$u = 0.09/\num{3.156e7} = \qty{2.85e-9}{m.s^{-1}}$ y $B = 2270\times\num{8e8}
\times\num{2.85e-9} = \qty{5.18e3}{kg.s^{-1}} = \qty{5.2}{Mg/s}$, unas
tres quintas partes de la de Hawái. El flujo de calor es $1170\times5178/\num{3.1e-5} =
\qty{1.95e11}{W}$, \qty{0.2}{TW}, o \qty{0.42}{\percent} de la pérdida de
calor de la Tierra.
\end{solution}

\begin{solution}{exo:b3:mantle-convection:8}
$v = 2\times30\times9.8\times(\num{4e5})^2/(9\times10^{21}) =
\qty{1.05e-8}{m.s^{-1}} = \qty{0.33}{m/yr}$; \qty{2000}{km} tardan
$\num{2e6}/\num{1.05e-8} = \qty{1.9e14}{s}$, unos \qty{6}{Myr}.
\end{solution}

\begin{solution}{exo:b3:mantle-convection:9}
$w = 300e^{-10/4.4} = \qty{30.9}{m}$ y $\dot w = w/\tau = 30.9/4400 =
\qty{7.0}{mm/yr}$, dos tercios de la velocidad medida cerca de Umeå: un
semiespacio uniforme con este $\tau$ está en el rango correcto, pero el rebote real es
más rápido o su depresión es más profunda.
\end{solution}

\begin{solution}{exo:b3:mantle-convection:10}
(a) Con $w = 2\rho_m\alpha\Delta T\sqrt{\kappa t/\pi}/(\rho_m - \rho_w)$,
\[
  \tfrac12(\rho_m - \rho_w)gw^2 = \frac{2g\rho_m^2\alpha^2\Delta T^2\kappa t}
  {\pi(\rho_m - \rho_w)} = g\rho_m\alpha\Delta T\kappa t\times
  \frac{2\rho_m\alpha\Delta T}{\pi(\rho_m - \rho_w)} .
\] (b) Por partes,
$\int_0^\infty x\operatorname{erfc}x\,\mathrm dx = [\tfrac{x^2}2\operatorname{erfc}
x]_0^\infty + \int_0^\infty\tfrac{x^2}2\,\tfrac2{\sqrt\pi}e^{-x^2}\,\mathrm dx =
0 + \tfrac1{\sqrt\pi}\cdot\tfrac{\sqrt\pi}4 = \tfrac14$. (c) A \qty{100}{Myr},
$\rho_mg\alpha\int z\theta = \qty{3.67e12}{N.m^{-1}}$ y, con $w =
\qty{3705}{m}$, $\tfrac12(\rho_m - \rho_w)gw^2 = \qty{1.53e11}{N.m^{-1}}$: el
término del agua es solo \qty{4}{\percent} de la fuerza.
\end{solution}

\begin{solution}{exo:b3:mantle-convection:11}
$\ln f = 3\ln(\pi^2 + k^2) - 2\ln k$, cuya derivada $6k/(\pi^2 + k^2) - 2/k$
se anula cuando $3k^2 = \pi^2 + k^2$, $k^2 = \pi^2/2$. Allí $f = (3\pi^2/2)^3/
(\pi^2/2) = 27\pi^4/4$. La longitud de onda $2\pi/k = 2\sqrt2\,d = 2.83\,d$ abarca un
par de celdas de rotación contraria, cada una con una anchura de $\sqrt2\,d = 1.41\,d$.
\end{solution}

\begin{solution}{exo:b3:mantle-convection:12}
(a) $u = \qty{1.90e-9}{m.s^{-1}}$, $\eta u/h = 10^{20}\times\num{1.90e-9}/
\num{1.5e5} = \qty{1.27}{MPa}$, y a lo largo de \qty{6000}{km}, \qty{7.6e12}{N.m^{-1}}.
(b) El doble del empuje de la dorsal y una sexta parte del tirón de la losa: el arrastre puede resistir
el empuje de la dorsal por sí solo, pero no la losa. (c) $\eta uL/h = \num{3.8e12} +
\num{4.6e12}$ da $u = \num{8.4e12}\times\num{1.5e5}/(\num{6e6}\times10^{20})
= \qty{2.1e-9}{m.s^{-1}} = \qty{6.6}{cm/yr}$, la velocidad de las placas oceánicas
rápidas.
\end{solution}

\begin{solution}{pb:b3:mantle-convection:1}
\textbf{1.} $p = 1000\times9.8\times300 = \qty{2.94}{MPa}$.

\textbf{2.} $w = \rho_wh/\rho_m = 300\times1000/3300 = \qty{90.9}{m}$. La corteza
se hunde como un bloque; lo que se desplaza desde debajo de la carga, y luego regresa, es el
manto, así que la columna se compensa con manto desplazado en profundidad.

\textbf{3.} $\pi\times(\num{1.25e5})^2 = \qty{4.91e10}{m^2}$, y el agua
pesa $\num{4.91e10}\times300\times1000 = \qty{1.47e16}{kg}$.

\textbf{4.} La línea de costa se formó nivelada, cuando el fondo estaba más deprimido
en el centro; el centro se ha elevado más que el borde desde entonces, así que la línea
ahora asciende hacia el centro, \qty{77.2}{m} a lo largo de \qty{125}{km}:
\qty{0.62}{m/km}.

\textbf{5.} $\lambda/2 = \qty{250}{km}$, luego $\lambda = \qty{500}{km}$ y $k =
2\pi/\num{5e5} = \qty{1.257e-5}{m^{-1}}$.

\textbf{6.} $1/k = \lambda/2\pi = \qty{79.6}{km}$.

\textbf{7.} $w = 90.9 - \text{levantamiento}$: 79.4, 69.3, 52.9, 40.4, 26.9 y
\qty{13.7}{m}.

\textbf{8.} $w = w_0e^{-t/\tau}$ da $\ln w = \ln w_0 - t/\tau$, una recta
de pendiente $-1/\tau$.

\textbf{9.} $\tau = (14 - 1)/\ln(79.4/13.7) = 13/1.7571 = \qty{7.40}{kyr}$.

\textbf{10.} $w(6) = 79.4\,e^{-5/7.40} = \qty{40.4}{m}$, con un levantamiento de
\qty{50.5}{m}, como en la tabla.

\textbf{11.} Quedan \qty{13.7}{m}; $\dot w = 13.7/7400 = \qty{1.85}{mm/yr}$.

\textbf{12.} Los \qty{13.7}{m} enteros, alcanzados cada vez más lentamente.

\textbf{13.} La velocidad disminuye como $e^{-t/\tau}$: $t = 7.40\ln(1.85/0.5) =
7.40\times1.309 = \qty{9.7}{kyr}$, unos 9700 años más.

\textbf{14.} $\tau = 7399\times\num{3.156e7} = \qty{2.335e11}{s}$, y $\eta =
\tau\rho_mg\lambda/(4\pi) = \num{2.335e11}\times3300\times9.8\times\num{5e5}/
12.566 = \qty{3.0e20}{Pa.s}$.

\textbf{15.} Justo por debajo del extremo inferior del intervalo fenoescandinavo. El flujo
bajo una carga de \qty{500}{km} de longitud de onda decae en \qty{80}{km}: el
rebote muestrea el manto más somero y débil, bajo y en parte dentro de lo que
en la Tierra real, es la litosfera.

\textbf{16.} $\tau = 4\pi\times\num{3.0e20}/(3300\times9.8\times\num{3e6}) =
\qty{3.9e10}{s} = \qty{1.2}{kyr}$. Bajo una carga tan amplia el flujo alcanza
hasta \qty{477}{km} y más, hacia un manto inferior diez veces más viscoso,
que lo frena.

\textbf{17.} Sostiene en parte una carga estrecha por flexión: la depresión en
equilibrio es más somera y más ancha que los \qty{90.9}{m} isostáticos, y parte
del rebote es elástica e inmediata. Un semiespacio viscoso uniforme ajustado al
levantamiento da entonces una viscosidad que mezcla ambos efectos.

\textbf{18.} $\tau_M = \num{3.0e20}/\num{7.7e10} = \qty{3.9e9}{s}$, unos
\qty{124}{yr}, sesenta veces más corto que $\tau$: a lo largo de miles de años el
manto responde como un fluido viscoso.

\textbf{19.} $1 - e^{-50/7.40} = 0.9988$: se alcanzó el \qty{99.9}{\percent} de la
depresión de equilibrio, y el supuesto se cumple.

\textbf{20.} $\eta \propto \lambda$: al reducir $\lambda$ a la mitad, $\eta$ se reduce a la mitad, y
$\log_{10}\eta$ disminuye en $\log_{10}2 = 0.30$.

\textbf{21.} $\mathrm{Ra} = 3700\times9.9\times\num{3.1e-5}\times1450\times
(\num{6.7e5})^3/(10^{-6}\times\num{3.0e20}) = \num{1.65e6}$.

\textbf{22.} $\num{1.65e6}/657.5 \approx 2500$ veces el valor crítico.

\textbf{23.} $(\mathrm{Ra}/2\mathrm{Ra}_c)^{1/3} = (1254)^{1/3} = 10.78$: $\delta
= 670/10.78 = \qty{62}{km}$ y $\mathrm{Nu} = 5.4$.

\textbf{24.} $u = (10^{-6}/\num{6.7e5})\times10.78^2 = \qty{1.73e-10}{m.s^{-1}}
= \qty{0.55}{cm/yr}$, de dos a catorce veces más lenta que las placas.

\textbf{25.} $q = 5.4\times3.1\times1450/\num{6.7e5} = \qty{36}{mW.m^{-2}}$. La
capa excluye el calor producido dentro del manto y el calor procedente del
núcleo y del manto inferior; y el escalamiento omite sus constantes de orden uno,
mientras que las placas, que se enfrían como semiespacios desde las dorsales, pierden mucho más calor sobre
suelo oceánico joven que el que perdería una capa límite media.

\textbf{26.} $\eta = \qty{3.0e20}{Pa.s}$: $\log_{10}\eta =
\textbf{\num[text-series-to-math=true]{20.48}}$.
\end{solution}
